\documentclass[10pt,a4paper]{amsart}
\usepackage{amsmath,amssymb,amsthm,graphicx}
% Retained preamble for the Stacks Project chapter {CH} (shared project preamble).
% The chapter file pinned-original-source.tex only carries \input{{preamble}} on its first line,
% so no preamble text lives inside the chapter itself.  The Stacks Project uses one single shared
% preamble for every chapter: the snapshot stages it three times (sources/stacks-chow/preamble.tex,
% sources/stacks-etale-cohomology/preamble.tex and focused-seeds/AGP/sources/stacks-snapshot-20261011/
% preamble.tex) and those copies are identical in size.  The block below is copied mechanically,
% line for line, from sources/stacks-chow/preamble.tex lines 14-18 and 154-217: the Xy-pic setup and
% the native theorem declarations together with the mathematical macros.  Nothing was added, removed
% or redefined; the wrapper class and the amsmath/amssymb/amsthm/graphicx base packages are supplied
% by the assembler.
% For commutative diagrams we use Xy-pic
% Chapter: stacks-cohomology.tex -- Cohomology of Sheaves.
\usepackage[all]{xy}

% We use 2cell for 2-commutative diagrams.
\xyoption{2cell}

% Theorem environments.
%
\theoremstyle{plain}
\newtheorem{theorem}[subsection]{Theorem}
\newtheorem{proposition}[subsection]{Proposition}
\newtheorem{lemma}[subsection]{Lemma}

\theoremstyle{definition}
\newtheorem{definition}[subsection]{Definition}
\newtheorem{example}[subsection]{Example}
\newtheorem{exercise}[subsection]{Exercise}
\newtheorem{situation}[subsection]{Situation}

\theoremstyle{remark}
\newtheorem{remark}[subsection]{Remark}
\newtheorem{remarks}[subsection]{Remarks}

\numberwithin{equation}{subsection}

% Macros
%
\def\lim{\mathop{\mathrm{lim}}\nolimits}
\def\colim{\mathop{\mathrm{colim}}\nolimits}
\def\Spec{\mathop{\mathrm{Spec}}}
\def\Hom{\mathop{\mathrm{Hom}}\nolimits}
\def\Ext{\mathop{\mathrm{Ext}}\nolimits}
\def\SheafHom{\mathop{\mathcal{H}\!\mathit{om}}\nolimits}
\def\SheafExt{\mathop{\mathcal{E}\!\mathit{xt}}\nolimits}
\def\Sch{\mathit{Sch}}
\def\Mor{\mathop{\mathrm{Mor}}\nolimits}
\def\Ob{\mathop{\mathrm{Ob}}\nolimits}
\def\Sh{\mathop{\mathit{Sh}}\nolimits}
\def\NL{\mathop{N\!L}\nolimits}
\def\CH{\mathop{\mathrm{CH}}\nolimits}
\def\proetale{{pro\text{-}\acute{e}tale}}
\def\etale{{\acute{e}tale}}
\def\QCoh{\mathit{QCoh}}
\def\Ker{\mathop{\mathrm{Ker}}}
\def\Im{\mathop{\mathrm{Im}}}
\def\Coker{\mathop{\mathrm{Coker}}}
\def\Coim{\mathop{\mathrm{Coim}}}

% Boxtimes
%
\DeclareMathSymbol{\boxtimes}{\mathbin}{AMSa}{"02}

%
% Macros for moduli stacks/spaces
%
\def\QCohstack{\mathcal{QC}\!\mathit{oh}}
\def\Cohstack{\mathcal{C}\!\mathit{oh}}
\def\Spacesstack{\mathcal{S}\!\mathit{paces}}
\def\Quotfunctor{\mathrm{Quot}}
\def\Hilbfunctor{\mathrm{Hilb}}
\def\Curvesstack{\mathcal{C}\!\mathit{urves}}
\def\Polarizedstack{\mathcal{P}\!\mathit{olarized}}
\def\Complexesstack{\mathcal{C}\!\mathit{omplexes}}
% \Pic is the operator that assigns to X its picard group, usage \Pic(X)
% \Picardstack_{X/B} denotes the Picard stack of X over B
% \Picardfunctor_{X/B} denotes the Picard functor of X over B
\def\Pic{\mathop{\mathrm{Pic}}\nolimits}
\def\Picardstack{\mathcal{P}\!\mathit{ic}}
\def\Picardfunctor{\mathrm{Pic}}
\def\Deformationcategory{\mathcal{D}\!\mathit{ef}}
\usepackage{cleveref}
\crefname{originalreference}{original reference}{original references}
\Crefname{originalreference}{Original reference}{Original references}
\makeatletter
\newlabel{situation-eta}{{0GT4}{source}{}{sourcepointer-1}{}}
\newlabel{situation-eta@cref}{{[originalreference][0][]0GT4}{[1][0][]source}{}{}{}}
\makeatother
\crefname{originalequation}{original equation}{original equations}
\Crefname{originalequation}{Original equation}{Original equations}
\makeatletter
\newlabel{equation-complex-bocksteins}{{0GT8}{source}{}{sourcepointer-2}{}}
\newlabel{equation-complex-bocksteins@cref}{{[originalequation][0][]0GT8}{[1][0][]source}{}{}{}}
\makeatother
\crefname{originallemma}{original lemma}{original lemmas}
\Crefname{originallemma}{Original lemma}{Original lemmas}
\makeatletter
\newlabel{lemma-eta-qis}{{0F8P}{source}{}{sourcepointer-3}{}}
\newlabel{lemma-eta-qis@cref}{{[originallemma][0][]0F8P}{[1][0][]source}{}{}{}}
\makeatother
\crefname{originalreference}{original reference}{original references}
\Crefname{originalreference}{Original reference}{Original references}
\makeatletter
\newlabel{more-algebra-lemma-eta-second-property}{{0F7T}{source}{}{sourcepointer-4}{}}
\newlabel{more-algebra-lemma-eta-second-property@cref}{{[originalreference][0][]0F7T}{[1][0][]source}{}{}{}}
\makeatother
\title{Canonical isomorphism between L-eta\_I(M) tensored with O\_X/I and the Bockstein complex H\textasciicircum{}bullet(M/I)}
\author{The Stacks Project Authors}
\date{stacks-project:cohomology.tex:local-snapshot-20261011; GNU FDL 1.2 or later}
\begin{document}
\maketitle
\noindent\textit{Editorial note.} All delivered mathematical text is the author's own, copied verbatim from the pinned Stacks Project chapter file cohomology.tex; no statement or proof step was rewritten, reordered or removed.\par\medskip
\noindent\textit{Editorial note.} The counted span cites four cross-chapter labels (including a situation and an equation), each recorded in assembly.external\_references as a navigation pointer with its canonical tag.\par\medskip
\noindent\textit{Editorial note.} The proof cites no bibliography item inside the counted span.\par\medskip
\noindent\textit{Allowed standard prerequisites (editorial).}\begin{itemize}
\item The functor L-eta\_I and its computation on torsion free complexes. Stacks Project, Cohomology of Sheaves, Lemma lemma-eta-qis (tag 0F8P) and Lemma lemma-Leta of this chapter. Justifies replacing M by a complex of I-torsion free modules and computing L-eta\_I there.
\item Algebraic eta second property. Stacks Project, More on Algebra, Lemma more-algebra-lemma-eta-second-property (tag 0F7T). The stalkwise model that the constructed map of complexes is compared against.
\end{itemize}\medskip
\ifdefined\hypertarget\hypertarget{sourcepointer-1}{}\fi
\noindent\textit{Source reference.} Original reference 0GT4: sources/stacks-chow/tags.txt, source lines 20954--20954. Cross-chapter standard prerequisite cited by the author inside the counted span; it is used as a pointer and its own proof is not part of this item.\par\medskip
\ifdefined\hypertarget\hypertarget{sourcepointer-2}{}\fi
\noindent\textit{Source reference.} Original equation 0GT8: sources/stacks-chow/tags.txt, source lines 20958--20958. Cross-chapter standard prerequisite cited by the author inside the counted span; it is used as a pointer and its own proof is not part of this item.\par\medskip
\ifdefined\hypertarget\hypertarget{sourcepointer-3}{}\fi
\noindent\textit{Source reference.} Original lemma 0F8P: sources/stacks-chow/tags.txt, source lines 19049--19049. Cross-chapter standard prerequisite cited by the author inside the counted span; it is used as a pointer and its own proof is not part of this item.\par\medskip
\ifdefined\hypertarget\hypertarget{sourcepointer-4}{}\fi
\noindent\textit{Source reference.} Original reference 0F7T: sources/stacks-chow/tags.txt, source lines 19018--19018. Cross-chapter standard prerequisite cited by the author inside the counted span; it is used as a pointer and its own proof is not part of this item.\par\medskip
\begin{lemma}
\label{lemma-eta-second-property}
In Situation \ref{situation-eta} let $M$ be an object of
$D(\mathcal{O}_X)$. There is a canonical isomorphism
$$
L\eta_\mathcal{I}M \otimes^\mathbf{L} \mathcal{O}_X/\mathcal{I}
\longrightarrow
H^\bullet(M/\mathcal{I})
$$
in $D(\mathcal{O}_X)$ where the right hand side is the complex
(\ref{equation-complex-bocksteins}).
\end{lemma}

\begin{proof}
By the construction of $L\eta_\mathcal{I}$ in Lemma \ref{lemma-eta-qis}
we may assume $M$ is represented by a complex of $\mathcal{I}$-torsion
free $\mathcal{O}_X$-modules $\mathcal{F}^\bullet$. Then
$L\eta_\mathcal{I}M$ is represented by the complex
$\eta_\mathcal{I}\mathcal{F}^\bullet$ which is a 
complex of $\mathcal{I}$-torsion free $\mathcal{O}_X$-modules as well.
Thus $L\eta_\mathcal{I}M \otimes^\mathbf{L} \mathcal{O}_X/\mathcal{I}$
is represented by the complex
$\eta_\mathcal{I}\mathcal{F}^\bullet \otimes \mathcal{O}_X/\mathcal{I}$.
Similarly, the complex $H^\bullet(M/\mathcal{I})$ has terms
$H^i(\mathcal{F}^\bullet \otimes \mathcal{I}^i/\mathcal{I}^{i + 1})$.

\medskip\noindent
Let $f$ be a local generator for $\mathcal{I}$.
Let $s$ be a local section of $(\eta_\mathcal{I}\mathcal{F})^i$.
Then we can write $s = f^is'$ for a local section $s'$ of
$\mathcal{F}^i$ and similarly $d^i(s) = f^{i + 1}t$ for a local
section $t$ of $\mathcal{F}^{i + 1}$. Thus $d^i$ maps $f^is'$
to zero in $\mathcal{F}^{i + 1} \otimes \mathcal{I}^i/\mathcal{I}^{i + 1}$.
Hence we may map $s$ to the class of $f^is'$ in
$H^i(\mathcal{F}^\bullet \otimes \mathcal{I}^i/\mathcal{I}^{i + 1})$.
This rule defines a map
$$
(\eta_\mathcal{I}\mathcal{F})^i \otimes \mathcal{O}_X/\mathcal{I}
\longrightarrow
H^i(\mathcal{F}^\bullet \otimes \mathcal{I}^i/\mathcal{I}^{i + 1})
$$
of $\mathcal{O}_X$-modules. A calculation shows that these maps
are compatible with differentials (essentially because $\beta$
sends the class of $f^is'$ to the class of $f^{i + 1}t$), whence
a map of complexes representing the arrow in the statement of the lemma.

\medskip\noindent
To finish the proof, we observe that the construction given
in the previous paragraph agrees on stalks with the maps
constructed in More on Algebra, Lemma
\ref{more-algebra-lemma-eta-second-property}
hence we conclude.
\end{proof}

\end{document}
